\chapter*{LỜI MỞ ĐẦU}
\addcontentsline{toc}{chapter}{Lời mở đầu}

Bài tập lớn của nhóm cài đặt thuật toán NSGA-II cho bài toán định tuyến xe xanh đa mục tiêu có khung
thời gian (MOGVRPTW), dựa trên mô hình của bài báo Li và cộng sự năm 2025~\cite{li2025}. Bài toán có ba mục tiêu xung
đột nhau: chi phí giao hàng, lượng khí thải CO\textsubscript{2} và mức hài lòng của khách hàng. Sau khi
có thuật toán gốc, nhóm bổ sung bốn cải tiến. Mỗi cải tiến do một thành viên phụ trách, có cờ bật tắt
riêng trong tệp cấu hình và mặc định ở trạng thái tắt, để có thể đo đóng góp của từng cải tiến.

Báo cáo này trình bày phần việc của thành viên thứ hai (TV2) gồm hai cải tiến: khởi tạo quần thể thông
minh (\code{init\_smart}) và tìm kiếm cục bộ trong từng tuyến (\code{local\_search}). Ngoài hai mô-đun
thuật toán, phần việc còn gồm bộ kiểm thử, gắn cải tiến vào vòng lặp NSGA-II của nhóm, kiểm tra hồi
quy để bảo đảm khi tắt cờ thì kết quả trùng khớp với thuật toán gốc, chạy thực nghiệm so sánh trên
nhiều bộ dữ liệu và vẽ các biểu đồ kết quả.

Nội dung được trình bày theo thứ tự: bài toán và nhiệm vụ của TV2 (Chương~1), cơ sở lý thuyết
(Chương~2), thiết kế và cài đặt khởi tạo thông minh (Chương~3), thiết kế và cài đặt tìm kiếm cục bộ
(Chương~4), tích hợp vào NSGA-II và kiểm thử (Chương~5), thực nghiệm (Chương~6), tiến độ cùng những
việc còn cần hoàn thiện (Chương~7). Mọi số liệu trong báo cáo được lấy từ các lần chạy thật trên mã
nguồn của nhóm; cách tái lập từng bảng và từng hình được ghi ở phần phụ lục.
